% ============================================================================
% Fichier  : partie5/P05_C23_decider.tex
% Titre    : Décider en Situation de Crise
% Auteur   : MINKA MI NGUIDJOÏ Thierry Emmanuel
% Version  : 1.0 -- novembre 2026
% ============================================================================
% Sources : Klein (1998) ; Kahneman (2011) ; Janis (1972) ; Weick (1993).
% Décision sur la rançon de l'Opération GAROUA : cas fictif.
% ============================================================================

\chapter{Décider en Situation de Crise}
\label{chap:decider}
\chaptermark{Décider en crise}

\epigraph{%
    « En crise, ne pas décider est aussi une décision,\\
    et souvent la pire. »%
}{Principe du cours}

\niveauChapitre{Expert}{Chapitres~\ref{chap:risque},
\ref{chap:perception} et~\ref{chap:cellule}.}

\begin{boxobjectifs}
\begin{itemize}
    \item Expliquer comment les experts décident sous pression (modèle de
          Klein).
    \item Reconnaître les biais propres à la crise~: effet tunnel, pensée de
          groupe, biais de normalité.
    \item Distinguer décisions réversibles et irréversibles, et les traiter
          différemment.
    \item Structurer une décision difficile~: options, critères, conséquences.
    \item Préparer à froid les décisions que l'on devra prendre à chaud.
\end{itemize}
\end{boxobjectifs}

\bigskip

% ============================================================================
\section*{Question directrice : comment décider de payer ou non une rançon en quarante-huit heures~?}
% ============================================================================

Le 7~mars, la cellule de crise doit trancher~: payer 85\,000~dollars pour
une clé de déchiffrement et la promesse de ne pas publier les données, ou
refuser. L'information est incomplète, le temps compte, l'enjeu est grand,
et la décision ne pourra pas être défaite. Ce chapitre donne une méthode
pour ce type de décision.

\medskip
\textbf{Plan du chapitre.}
\begin{enumerate}
    \item Comment décident les experts (section~\ref{sec:c23-klein}).
    \item Les biais de la crise (section~\ref{sec:c23-biais}).
    \item Réversible ou irréversible (section~\ref{sec:c23-reversible}).
    \item Structurer une décision (section~\ref{sec:c23-structurer}).
    \item Décider à froid ce qui se décidera à chaud (section~\ref{sec:c23-froid}).
\end{enumerate}

\separateur

% ============================================================================
\section{Comment Décident les Experts}
\label{sec:c23-klein}
% ============================================================================

Gary Klein a observé des pompiers, des militaires, des médecins
urgentistes~: sous pression, ils ne comparent pas des options~; ils
\emph{reconnaissent} une situation à partir de leur expérience, imaginent
la première action plausible, la simulent mentalement et l'appliquent si
elle tient~\cite{klein1998}. C'est la \termecle{décision fondée sur la
reconnaissance}.

\begin{boxinfo}
\textbf{Conséquence pour la gestion de crise.} Ce mode de décision est
rapide et souvent bon, \emph{à condition} d'avoir déjà vu des situations
semblables. D'où l'importance des exercices (Chapitre~\ref{chap:exercices-crise})~:
ils donnent aux décideurs l'expérience que les crises réelles, rares, ne
leur donnent pas. Face à une situation vraiment nouvelle, il faut revenir
à une décision délibérée.
\end{boxinfo}

% ============================================================================
\section{Les Biais de la Crise}
\label{sec:c23-biais}
% ============================================================================

\begin{tableaucours}{L{3.4cm}L{4.6cm}L{4.2cm}}{%
    \tableauHeader{Biais} &
    \tableauHeader{Effet} &
    \tableauHeader{Parade}
}
Effet tunnel & Toute l'attention sur un seul problème (restaurer) &
Revue systématique des parties prenantes à chaque point \\
Biais de normalité & Croire au retour rapide à la normale & Hypothèse de
travail pessimiste, révisée à chaque point \\
Pensée de groupe & Unanimité de façade, objections tues & Tour de table
inversé, contradicteur désigné \\
Excès de confiance & Annoncer un délai de reprise trop court & Annoncer des
fourchettes, avec leur hypothèse \\
Fatigue & Décisions dégradées après de longues heures & Relèves, sommeil
imposé au directeur de crise \\
\end{tableaucours}
\vspace{-0.8em}
\begin{center}{\small\itshape Tableau~23.1. Biais de la crise et parades}\end{center}

\begin{boxscenario}{Opération GAROUA. « Lundi, tout sera rétabli »}
\textbf{Constatation.} Le 6~mars à 16~h, M.~KOUNOU annonce aux clients,
par la voix de M.~NGUELE, un retour à la normale « lundi 9~mars ». La
facturation reprend le 16~mars.

\medskip
\textbf{Analyse} (confiance élevée). L'annonce cumule biais de normalité et
excès de confiance~: on ignorait encore l'état des sauvegardes. Elle a
coûté de la crédibilité à chaque report. Une fourchette (« au plus tôt
lundi, au plus tard dans deux semaines, nous préciserons chaque jour »)
aurait été plus juste.
\end{boxscenario}

% ============================================================================
\section{Réversible ou Irréversible}
\label{sec:c23-reversible}
% ============================================================================

\begin{tableaucours}{L{3.2cm}L{4.4cm}L{4.6cm}}{%
    \tableauHeader{Décision} &
    \tableauHeader{Exemples} &
    \tableauHeader{Traitement}
}
Réversible & Couper les liaisons, isoler un serveur, basculer en mode
dégradé & Décider vite, sur l'information disponible~; corriger ensuite \\
Irréversible & Payer la rançon, publier un communiqué, licencier un
salarié, effacer un serveur & Prendre le temps strictement nécessaire~;
avis juridique~; options écrites \\
\end{tableaucours}
\vspace{-0.8em}
\begin{center}{\small\itshape Tableau~23.2. Deux catégories de décisions}\end{center}

\begin{boxalert}
\textbf{Violation typique.} Traiter avec la même lenteur la coupure des
liaisons (réversible) et le paiement de la rançon (irréversible).

\smallskip
\textbf{Règle~:} les décisions réversibles et conservatoires se prennent
tout de suite, au besoin par délégation~; les irréversibles se préparent.
Le 6~mars, la coupure a attendu 3~h~40 alors qu'elle se défaisait en une
minute.
\end{boxalert}

% ============================================================================
\section{Structurer une Décision}
\label{sec:c23-structurer}
% ============================================================================

\begin{boxprocedure}
\textbf{Structurer une décision irréversible}
\begin{enumerate}
    \item \textbf{Formuler la question} et l'échéance réelle.
    \item \textbf{Lister les options}, y compris des options intermédiaires.
    \item \textbf{Fixer les critères}~: continuité, coût, droit, éthique,
          réputation, précédent.
    \item \textbf{Évaluer} chaque option sur chaque critère, avec ce que
          l'on sait et ce que l'on ignore.
    \item \textbf{Décider}, consigner la décision, ses motifs et
          l'information disponible (Chapitre~\ref{chap:risque}).
\end{enumerate}
\end{boxprocedure}

\begin{tableaucours}{L{2.8cm}L{4.6cm}L{4.8cm}}{%
    \tableauHeader{Critère} &
    \tableauHeader{Payer} &
    \tableauHeader{Ne pas payer}
}
Continuité & Clé peut-être fournie~; aucune garantie qu'elle fonctionne &
Reprise par sauvegardes~: copie de fin janvier, ressaisie \\
Données & Promesse de non-publication sans garantie & Publication probable
d'un échantillon \\
Droit et éthique & Financement d'une activité criminelle~; sanctions
possibles selon le destinataire & Conforme à l'avis juridique \\
Précédent & Signal donné aux attaquants & Signal de fermeté \\
Coût & 85\,000~dollars et risque de seconde demande & Coût de la reprise
(Tableau~0.4) \\
\end{tableaucours}
\vspace{-0.8em}
\begin{center}{\small\itshape Tableau~23.3. La décision du 8 mars, structurée}\end{center}

\begin{boxscenario}{Opération GAROUA. La décision consignée}
\textbf{Main courante, 8~mars, 18~h~10.} « Décision~: LiZbiZ ne paie pas
la rançon. Motifs~: aucune garantie de déchiffrement ni de non-publication~;
avis de Me~TSALA~; sauvegarde de fin janvier exploitable selon
M.~KOUNOU. Information manquante~: étendue des données copiées. Décidé
par M.~ETEME. »

\medskip
\textbf{Analyse} (confiance élevée). La décision est documentée avec ses
motifs et ses inconnues~: elle pourra être jugée sur son dossier, et non
sur la publication du 19~mars.
\end{boxscenario}

% ============================================================================
\section{Décider à Froid ce qui se Décidera à Chaud}
\label{sec:c23-froid}
% ============================================================================

Les décisions les plus difficiles se préparent avant la crise~: on fixe à
froid une \emph{position de principe} et les \emph{délégations}.

\begin{tableaucours}{L{5.6cm}L{6.6cm}}{%
    \tableauHeader{Décision préparée} &
    \tableauHeader{Position adoptée par la direction de LiZbiZ (juin)}
}
Couper les liaisons entre sites & Délégation au responsable informatique
ou à l'astreinte, sans attendre le DG \\
Basculer en mode dégradé & Délégation au directeur des opérations \\
Payer une rançon & Principe~: non~; dérogation par le DG seul, après avis
juridique écrit \\
Premier message public & Message d'attente préapprouvé, publiable par la
communication \\
\end{tableaucours}
\vspace{-0.8em}
\begin{center}{\small\itshape Tableau~23.4. Décisions préparées à froid}\end{center}

\begin{boxmetier}{Directeur de crise}
\textbf{Déléguer les gestes conservatoires.} Le directeur de crise donne à
l'avance, par écrit, le pouvoir de prendre les décisions réversibles et
urgentes. Il se garde les irréversibles.
\end{boxmetier}

\separateur

\begin{boxresume}
\textbf{Ce qu'il faut retenir de ce chapitre~:}
\begin{itemize}
    \item \textbf{Sous pression}, les experts reconnaissent la situation~;
          les exercices donnent cette expérience.
    \item \textbf{Les biais de crise}~: tunnel, normalité, groupe,
          confiance, fatigue~; chacun a sa parade.
    \item \textbf{Réversible}~: vite~; \textbf{irréversible}~: préparé,
          documenté.
    \item \textbf{Une décision structurée}~: question, options, critères,
          évaluation, décision consignée.
    \item \textbf{À froid}, on fixe les positions de principe et les
          délégations.
\end{itemize}
\end{boxresume}

\begin{boxinfo}
\textbf{Points de contrôle de l'élève.} Je sais~: décrire la décision
fondée sur la reconnaissance~; reconnaître un biais de crise et sa
parade~; classer une décision en réversible ou irréversible~; structurer
une décision~; préparer des délégations.
\end{boxinfo}

\bigskip

\begin{boxexo}[title={\ding{45}\quad Exercice~23.1. Réversible \niveaubase}]
Classez en réversible ou irréversible~: isoler un poste~; envoyer un
communiqué de presse~; effacer un serveur pour le réinstaller~; basculer
sur Yaoundé.
\end{boxexo}

\begin{boxexo}[title={\ding{45}\quad Exercice~23.2. Les biais \niveaubase}]
Quel biais illustre l'annonce « lundi, tout sera rétabli »~? Quelle parade~?
\end{boxexo}

\begin{boxexo}[title={\ding{45}\quad Exercice~23.3. Klein \niveaubase}]
Pourquoi le modèle de Klein justifie-t-il les exercices de crise~?
\end{boxexo}

\begin{boxexo}[title={\ding{45}\quad Exercice~23.4. Une fourchette \niveaumoyen}]
Rédigez l'annonce de délai de reprise que LiZbiZ aurait dû faire le
6~mars à 16~h.
\end{boxexo}

\begin{boxexo}[title={\ding{45}\quad Exercice~23.5. Options intermédiaires \niveaumoyen}]
Entre « payer » et « ne pas payer », quelles options intermédiaires la
cellule pouvait-elle examiner~? Avec quelles limites~?
\end{boxexo}

\begin{boxexo}[title={\ding{45}\quad Exercice~23.6. Consigner \niveaumoyen}]
Pourquoi consigner l'information manquante au moment de la décision~?
\end{boxexo}

\begin{boxexo}[title={\ding{45}\quad Exercice~23.7. Délégations \niveaumoyen}]
Proposez deux autres décisions à déléguer à froid, avec leur délégataire.
\end{boxexo}

\begin{boxexo}[title={\ding{45}\quad Exercice~23.8. La fatigue \niveaumoyen}]
La cellule a siégé dix-huit heures le 7~mars. Proposez une organisation
des relèves pour le directeur de crise.
\end{boxexo}

\begin{boxexo}[title={\ding{45}\quad Exercice~23.9. Une décision bancaire \niveauexpert}]
Le système central de la BANQUE SANAGA est indisponible depuis trois
heures. Structurez la décision « fermer ou non les agences demain »~:
options, critères, évaluation.
\end{boxexo}

\begin{boxexo}[title={\ding{45}\quad Exercice~23.10. La politique \niveauexpert}]
Rédigez la politique de LiZbiZ sur les rançons~: principe, exceptions,
procédure de décision, documentation.
\end{boxexo}

\bigskip

\begin{boxinfo}
\textbf{Pour aller plus loin.}
\begin{itemize}
    \item G.~Klein, \emph{Sources of Power}~\cite{klein1998}.
    \item D.~Kahneman~\cite{kahneman2011}~; I.~Janis~\cite{janis1972}.
\end{itemize}
\end{boxinfo}

\begin{boxinfo}
\textbf{Transition.} Une décision doit être expliquée. Le
Chapitre~\ref{chap:communiquer} traite de la communication de crise.
\end{boxinfo}
